\documentclass{article}

% Packages
\usepackage{fancyhdr}
\usepackage{extramarks}
\usepackage{amsmath}
\usepackage{amsthm}
\usepackage{amsfonts}
\usepackage{tikz}
\usepackage[plain]{algorithm}
\usepackage{algpseudocode}
\usepackage{enumerate}
\usepackage{dsfont}
\usepackage{bbm}

\usetikzlibrary{automata,positioning}

% Document Layout
\topmargin=-0.45in
\evensidemargin=0in
\oddsidemargin=0in
\textwidth=6.5in
\textheight=9.0in
\headsep=0.25in
\linespread{1.1}

% Page Style
\pagestyle{fancy}
\lhead{\hmwkAuthorName}
\chead{\hmwkClass:\ \hmwkTitle}
\rhead{Section \hmwkSection, \firstxmark}
\lfoot{\lastxmark}
\cfoot{\thepage}
\renewcommand\headrulewidth{0.4pt}
\renewcommand\footrulewidth{0.4pt}

% Paragraph Settings
\setlength\parindent{0pt}
\setlength{\parskip}{5pt}

% Section Management
\newcommand{\hmwkSection}{A} % Current section (A, B, or C) - update manually

% Problem Header Management
\newcommand{\enterProblemHeader}[1]{
  \nobreak\extramarks{}{Problem \arabic{#1} continued on next page\ldots}\nobreak{}
  \nobreak\extramarks{Problem \arabic{#1} (continued)}{Problem \arabic{#1} continued on next page\ldots}\nobreak{}
}

\newcommand{\exitProblemHeader}[1]{
  \nobreak\extramarks{Problem \arabic{#1} (continued)}{Problem \arabic{#1} continued on next page\ldots}\nobreak{}
  \stepcounter{#1}
  \nobreak\extramarks{Problem \arabic{#1}}{}\nobreak{}
}

% Counters
\setcounter{secnumdepth}{0}
\newcounter{partCounter}
\newcounter{homeworkProblemCounter}
\setcounter{homeworkProblemCounter}{1}
\nobreak\extramarks{Problem \arabic{homeworkProblemCounter}}{}\nobreak{}

% Homework Problem Environment
% Optional argument adjusts problem counter for non-sequential problems
\newenvironment{homeworkProblem}[1][-1]{
  \ifnum#1>0
    \setcounter{homeworkProblemCounter}{#1}
  \fi
  \section{Problem \arabic{homeworkProblemCounter}}
  \setcounter{partCounter}{1}
  \enterProblemHeader{homeworkProblemCounter}
}{
  \exitProblemHeader{homeworkProblemCounter}
}

% Assignment Details
\newcommand{\hmwkTitle}{Sheet\ \#1}
\newcommand{\hmwkDueDate}{Feb 1, 2026}
\newcommand{\hmwkClass}{C6.2 Continuous Optimization}
\newcommand{\hmwkClassInstructor}{Professor C. Cartis}
\newcommand{\hmwkAuthorName}{\textbf{Ray Tsai}}

% Title Page
\title{
  \vspace{2in}
  \textmd{\textbf{\hmwkClass:\ \hmwkTitle}}\\
  \normalsize\vspace{0.1in}\small{Due\ on\ \hmwkDueDate\ at 12:00pm}\\
  \vspace{0.1in}\large{\textit{\hmwkClassInstructor}} \\
  \vspace{3in}
}

\author{\hmwkAuthorName}
\date{}

% Part Command
\renewcommand{\part}[1]{\textbf{\large Part \Alph{partCounter}}\stepcounter{partCounter}\\}

% Mathematical Commands
% Algorithms
\newcommand{\alg}[1]{\textsc{\bfseries \footnotesize #1}}

% Calculus
\newcommand{\deriv}[1]{\frac{\mathrm{d}}{\mathrm{d}x} (#1)}
\newcommand{\pderiv}[2]{\frac{\partial}{\partial #1} (#2)}
\newcommand{\dx}{\mathrm{d}x}

% Probability and Statistics
\newcommand{\Var}{\mathrm{Var}}
\newcommand{\Cov}{\mathrm{Cov}}
\newcommand{\Bias}{\mathrm{Bias}}
\newcommand*{\prob}{\mathds{P}}
\newcommand*{\E}{\mathds{E}}

% Number Sets
\newcommand*{\Z}{\mathbb{Z}}
\newcommand*{\Q}{\mathbb{Q}}
\newcommand*{\R}{\mathbb{R}}
\newcommand*{\C}{\mathbb{C}}
\newcommand*{\N}{\mathbb{N}}

\begin{document}

\maketitle
\pagebreak

\begin{homeworkProblem}
  Consider the function
  \[ 
    f(x) = 10(x_2 - x_1^2)^2 + (1 - x_1)^2, \quad x = (x_1 \quad x_2)^T \in \mathbb{R}^2. 
  \]

  \begin{enumerate}[(a)]
    \item Compute the gradient vector and the Hessian matrix of $f$ at (any) $x \in \mathbb{R}^2$. Find all stationary points of $f$. Show that $x^* = (1 \quad 1)^T$ is the unique global minimizer of $f$ and that the Hessian of $f$ at $x^*$ is positive definite.

    \begin{proof}
      \[
        \nabla f(x) = \begin{bmatrix}
          -40x_1(x_2 - x_1^2) - 2(1 - x_1) \\
          20(x_2 - x_1^2)
        \end{bmatrix}, \quad
        \nabla^2 f(x) = \begin{bmatrix}
          -40(x_2 - x_1^2) + 80x_1^2 + 2 & -40x_1 \\
          -40x_1 & 20
        \end{bmatrix}.
      \]
      Note that $\nabla f(x) = 0$ only if $x = (1, 1)^T$. Since all minors of $\nabla^2 f(1, 1) = \begin{bmatrix}
        82 & -40 \\
        -40 & 20
      \end{bmatrix}$ are positive, $\nabla^2 f$ is positive definite at $(1, 1)^T$.
    \end{proof}
    
    \item Show that the Hessian matrix $\nabla^2 f(x)$ of $f$ is singular if and only if $x$ satisfies the condition
    \[ x_2 - x_1^2 = 0.05. \]
    Hence show that $\nabla^2 f(x)$ is positive definite for all $x$ such that $f(x) < 0.025$.

    \begin{proof}
      Note that
      \[
        \det(\nabla^2 f(x)) = 800(x_1^2 - x_2) + 40.
      \]
      Thus $\nabla^2 f(x)$ is singular if and only if $x_2 - x_1^2 = 40/800 = 0.05$. Thus if $f(x) < 0.025$, then 
      \[
        10(x_2 - x_1^2)^2 \leq 10(x_2 - x_1^2)^2 + (1 - x_1)^2 < 0.025 \implies x_2 - x_1^2 < 0.05.
      \]
    \end{proof}
    
    \item Show that $f$ is not a convex function.
    \begin{proof}
      Note that 
      \[
        \nabla^2 f(0, 1) = \begin{bmatrix}
          -38 & 0 \\
          0 & 20
        \end{bmatrix}
      \]
      has negative determinant, so $\nabla^2 f(x)$ is not semipositive definite for all $x$. Thus, $f$ is not a convex function.
    \end{proof}
  \end{enumerate}
\end{homeworkProblem}

\newpage

\begin{homeworkProblem}
  Show that the function
  \[ 
    f(x) = (x_2 - x_1^2)^2 + x_1^5
  \]
  has only one stationary point which is neither a local maximum nor a local minimum.

  \begin{proof}
    \[
      \nabla f(x) = \begin{bmatrix}
        5x_1^4 + 4x^3_1 - 4x_1x_2 \\
        2x_2 - 2x_1^2
      \end{bmatrix}, \quad
      \nabla^2 f(x) = \begin{bmatrix}
        20x_1^3 + 12x_1^2 - 4x_2 & -4x_1 \\
        -4x_1 & 2
      \end{bmatrix}.
    \]
    Note that $\nabla f(x) = 0$ only when $x = (0, 0)^T$. Suppose $x_2 = x_1^2$. Then $f(x) = x_1^5$, which has a saddle point at $(0, 0)^T$. This shows that $x = (0, 0)^T$ is a stationary point which is neither a local maximum nor a local minimum.
  \end{proof}
\end{homeworkProblem}

\newpage

% To change sections, use: \renewcommand{\hmwkSection}{B}
% Example: Uncomment the line below when you start Section B
\renewcommand{\hmwkSection}{B}

\begin{homeworkProblem}
  question

  \begin{proof}
    proof
  \end{proof}
\end{homeworkProblem}

\end{document}